\section{Conclusions}
\label{sec:conclusion}

This paper presented a calibrated, modular geospatial optimization framework
for retail branch expansion that integrates pedestrian-network accessibility,
graph-based spectral clustering, MCA-based socioeconomic demand weighting,
and a clustered \mbox{p-median} formulation.
Relative to the original submission, the framework incorporates three
substantive advances: parameter calibration from a Pareto-frontier grid search,
a demand-proportional cluster-budget allocation rule, and a rigorous baseline
comparison conducted entirely under pedestrian-network distance.

The empirical results across three urban contexts in Mexico establish two
principal findings.
First, territorial decomposition via spectral clustering retains
93--100\% of the global \mbox{p-median} coverage while reducing
solver runtime by up to 67\% in medium-to-large instances;
the decomposition gap is largest in supply-constrained peri-urban
settings where most blocks qualify as candidates.
Second, Euclidean-distance models systematically underestimate coverage
and overestimate accessibility by 10--20\%, confirming that network-based
distances are not a refinement but a requirement for urban facility location.

The methodological contribution is the demonstration that MCA and spectral
clustering can be integrated into a single, calibrated facility-location
pipeline whose parameters are selected from data rather than convention,
and whose performance can be benchmarked against a provably optimal global
reference.
The pipeline is modular: each component can be replaced without
redesigning the whole system, and the parameter calibration procedure
transfers directly to new study areas.

Future work can extend the framework in several directions.
Incorporating site acquisition and operating costs would make it directly
actionable for investment decisions.
Competition-sensitive demand models---where each new opening also affects
the patronage of nearby existing stores---would extend the commercial
validity of the objective.
Scenario analysis on service-radius assumptions, robustness under uncertain
travel times, and validation against observed transaction data would further
strengthen the framework as a decision-support tool for retail and
public-service network design.
